\documentclass[11pt]{article}
\usepackage[letterpaper,margin=1in]{geometry}
\usepackage{amsmath,amssymb,booktabs,hyperref}
\hypersetup{hidelinks,pdftitle={Certified Fourier bounds for the least quadratic nonresidue}}
\newcommand{\paperAuthor}{Lennart Lopin\\ {\normalsize Euler's Identity LLC}\\ {\normalsize ORCID: \href{https://orcid.org/0009-0001-2751-2060}{0009-0001-2751-2060}}}
% Enter the public release DOI URL or a permanent repository/commit URL here.
% Filling this one field enables the in-text software citation and bibliography item.
\newcommand*{\replicationURL}{https://doi.org/10.5281/zenodo.24408857}
\newcommand*{\ifReplicationURL}[1]{\ifx\replicationURL\empty\else#1\fi}
\title{Certified Fourier bounds for the least quadratic nonresidue}
\author{\paperAuthor}
\date{November 20, 2026}
\newcommand{\CC}{\mathcal C(1)}
\newcommand{\abs}[1]{\left|#1\right|}
\newcommand{\RR}{\mathbb R}
\newcommand{\ii}{\mathrm i}
\begin{document}
\maketitle
\begin{abstract}
We give two explicit Fourier test functions for the extremal constant
$\mathcal C(1)$ of Carneiro, Milinovich, Quesada-Herrera and Ramos. A
single-shift function with 36 rational coefficients certifies
$\mathcal C(1)^{-2}<0.760777412$. A two-shift function with 24 rational
coefficients certifies $\mathcal C(1)^{-2}<0.76079030$. Their published
upper bound is $0.7615$. At least two separately written interval
computations check each witness. Under the generalized Riemann
hypothesis for Dirichlet $L$-functions, their Theorems 1 and 2 give the
bound $0.760777412$ for the limiting ratios of the least quadratic
nonresidue and the least prime quadratic residue to $(\log p)^2$.
The Fourier inequalities themselves are unconditional.
\end{abstract}

\section{The bounds and their consequences}
For an odd prime $p$, let $n_p$ be the least positive quadratic nonresidue
modulo $p$ and $r_p$ the least prime that is a nonzero quadratic residue
modulo $p$. Write $\widehat F(t)=\int_{\RR} F(x)e^{-2\pi\ii xt}\,dx$. For
nonzero $F\in L^1(\RR)$ with real-valued transform, define $D=\|F\|_1$ and
\[
 N=2\pi\left(\int_{-\infty}^0 \widehat F(t)e^{\pi t}\,dt
                 -\int_0^\infty |\widehat F(t)|e^{\pi t}\,dt\right).
\]
The supremum of $N/D$ is $\CC$ in \cite[(1.1)]{cmqr}. Each admissible
witness with $N>0$ therefore gives $\CC^{-2}\le(D/N)^2$.

\medskip\noindent
\textbf{Proposition.} The two-shift function of Section 2 satisfies
$(D/N)^2<0.76079030$, and the single-shift function of Section 3
satisfies $(D/N)^2<0.760777412$. Consequently, assuming Dirichlet GRH,
\[
 \max\!\left\{\limsup_{p\to\infty}\frac{n_p}{(\log p)^2},
               \limsup_{p\to\infty}\frac{r_p}{(\log p)^2}\right\}
 \le \CC^{-2}<0.760777412 ,
\]
where $p$ runs over primes. The first inequality is precisely the joint
consequence of Theorems 1 and 2 of \cite{cmqr}. The improvement is in the
coefficient of $(\log p)^2$.

Equivalently, $\CC>1.1464924$, compared with $\CC>1.14600$ in
\cite[Theorem 5]{cmqr}. The same theorem gives $\CC<1.14731$, so
$\CC^{-2}>0.7596$ and the method cannot go below that value. The
single-shift witness removes about $37.6\%$ of the interval for $\CC^{-2}$
left open by \cite[Theorem 5]{cmqr}, and the two-shift witness alone
removes about $36.9\%$. The two witnesses belong to different families.

Alp\"oge's announcement of a Lean formalization covering Dirichlet
GRH~\cite{alpoge} motivated the search for these witnesses. Sutherland reports
independent checks of the formal statements~\cite{sutherland}. Taking the
Dirichlet GRH result as established discharges the hypothesis of both
arithmetic corollaries above.

\newpage
\section{A two-shift witness}
Take $a=0.246$, $c_1=0.85$, $c_2=0$, and use the integers in Table 1 to
define $b_{jn}=B_{jn}/10^7$. Put
\[
 P_j(u)=\sum_{n=1}^{12} b_{jn}\frac{u^{2n-1}}{(2n-1)!},\qquad
 u_j=\frac{c_j-\pi t}{a},\qquad
 \widehat F(t)=\sum_{j:\,u_j\ge0}e^{-u_j}P_j(u_j).
\]
Fourier inversion and the gamma integral give
\[
 F(x)=\frac a\pi\sum_{j,n} b_{jn} e^{2\ii c_jx}(1+2\ii ax)^{-2n}.
\]
Thus $F\in L^1$, and its Fourier transform is real and continuous. The
two translations allow the shape on negative frequencies to change while
only the first component reaches positive frequencies.

\begin{table}[h]
\centering
\begin{tabular}{r rr@{\qquad}r rr}
\toprule
$n$ & $B_{1n}$ & $B_{2n}$ & $n$ & $B_{1n}$ & $B_{2n}$ \\
\midrule
1 & 1295 & 0 & 7 & -46679113 & -118505260 \\
2 & 79929 & -9074 & 8 & 20002165 & 154829818 \\
3 & -6609595 & 5043626 & 9 & -2996805 & -138162493 \\
4 & 117992958 & -14383101 & 10 & -22963877 & 320107061 \\
5 & -234402020 & -118792841 & 11 & 279853645 & 234086697 \\
6 & 332397224 & 61384470 & 12 & -809901785 & 62986550 \\
\bottomrule
\end{tabular}
\caption{Two-shift witness: exact integer coefficients in ascending odd powers.}
\end{table}

\paragraph{Numerator.} Set $w=\pi t$ and $\lambda=1+a$. Integrating the
signed transform over all $w$ and subtracting twice its positive part on
$w>0$ yields
\[
 N=2a\sum_{j,n}\frac{e^{c_j}b_{jn}}{(1+a)^{2n}}
       -4ae^{c_1}\int_0^{c_1/a}e^{-\lambda u}\max\{P_1(u),0\}\,du .
\]
Here $\int_0^\infty e^{-\lambda u}u^{2n-1}/(2n-1)!\,du=\lambda^{-2n}$.

\paragraph{Norm.} After removing a common phase, let $z=2ax$,
$\omega_j=(c_j-c_1)/a$, and put
\[
 H(z)=\sum_{j,n} b_{jn} e^{\ii\omega_j z}(1+\ii z)^{-2n},\qquad
 D=\frac1\pi\int_0^\infty |H(z)|\,dz .
\]
The factor $1/\pi$ uses $H(-z)=\overline{H(z)}$ for real $z$. For $R>0$,
the triangle inequality and $|1+\ii z|\ge z$ for $z>0$ give
\begin{equation}
 \int_R^\infty |H(z)|\,dz
 \le\sum_{j,n}\frac{|b_{jn}|}{(2n-1)R^{2n-1}} . \label{eq:tail}
\end{equation}

\newpage
\section{A single-shift witness with 36 coefficients}
In this section take $a=3/10$ and $c=4/5$, and let
$S(z)=\sum_{j=0}^{35}r_jz^j$ with $r_j=q_j/10^6$ and the integers $q_j$
of Table 2. (The verifiers call this polynomial $R$.) Define $d_k$ and
$P$ by
\[
 S(2w-1)=\sum_{k=1}^{36}d_k w^{k-1},
 \qquad P(u)=\sum_{k=1}^{36}d_k\,\frac{u^k}{k!}.
\]
Take $\widehat F(t)=\mathbf 1_{\{u\ge0\}}e^{-u}P(u)$, where
$u=(c-\pi t)/a$. This transform is real, continuous at the cutoff because
$P(0)=0$, and integrable. Substitution into the inverse Fourier integral
gives
\[
 F(x)=\frac a\pi e^{2\ii cx}\sum_{k=1}^{36}
      \frac{d_k}{k!}\int_0^\infty e^{-(1+2\ii ax)u}u^k\,du
     =\frac a\pi e^{2\ii cx}\sum_{k=1}^{36}\frac{d_k}{(1+2\ii ax)^{k+1}}.
\]
Thus $F(x)=O((1+|x|)^{-2})$, Fourier inversion recovers the prescribed
continuous transform, and $S\ne0$ ensures $F\ne0$.

\begin{table}[h]
\centering
\setlength{\tabcolsep}{8pt}
\begin{tabular}{@{}rr@{\qquad}rr@{\qquad}rr@{\qquad}rr@{}}
\toprule
$j$&$q_j$&$j$&$q_j$&$j$&$q_j$&$j$&$q_j$\\
\midrule
0 & 1000000 &  9 & 510409 & 18 & 6332    & 27 & $-1653$\\
1 & 4815715 & 10 & 258500 & 19 & 5339    & 28 & 1541\\
2 & 8710552 & 11 & 235574 & 20 & $-4625$ & 29 & $-1050$\\
3 & 7617483 & 12 & 72629  & 21 & 6832    & 30 & 1019\\
4 & 4487466 & 13 & 71024  & 22 & $-4287$ & 31 & $-464$\\
5 & 3434816 & 14 & 46027  & 23 & 2822    & 32 & 1320\\
6 & 2345152 & 15 & 3286   & 24 & $-479$  & 33 & 215\\
7 & 1189056 & 16 & 28100  & 25 & $-732$  & 34 & 2070\\
8 & 938349  & 17 & $-2855$ & 26 & 1601   & 35 & 5064\\
\bottomrule
\end{tabular}
\caption{Single-shift witness: $r_j=q_j/10^6$ is the coefficient of $z^j$ in $S(z)$.}
\end{table}

\paragraph{Norm.} By evenness of $|F|$, put $\theta=\arctan(2ax)$ for
$x\ge0$. Writing $w=(1+\ii\tan\theta)^{-1}=\cos\theta\,e^{-\ii\theta}$ gives
$2w-1=e^{-2\ii\theta}$. The factor $w^2$ in the sum for $F$ has modulus
$\cos^2\theta$, which cancels the $\sec^2\theta$ in $2a\,dx$. Hence
\begin{equation}
 D_*:=\pi\|F\|_1=\int_0^{\pi/2}|S(e^{-2\ii\theta})|\,d\theta. \label{eq:normS}
\end{equation}

\paragraph{Numerator.} Put $h=1+a=13/10$ and $m=c/a=8/3$. Here $t<0$
corresponds to $u>m$, and $0<t<c/\pi$ corresponds to $0<u<m$. Using
$e^{\pi t}=e^{c-au}$, the quantity $N_*:=\pi N$ satisfies
\begin{equation}
 N_*=2\pi a e^c\left(\int_m^\infty e^{-hu}P(u)\,du-
                    \int_0^m e^{-hu}|P(u)|\,du\right), \label{eq:numS}
\end{equation}
and $N/D=N_*/D_*$. For $v\ge0$, repeated integration by parts gives
\[
 T(v)=\int_v^\infty e^{-hu}P(u)\,du
     =e^{-hv}\sum_{j=0}^{36}\frac{v^j}{j!}
                    \sum_{k=j}^{36}\frac{d_k}{h^{k-j+1}},
 \qquad d_0=0.
\]
Exact rational real-root isolation shows that $P$ has a simple root at
zero and exactly two simple roots in $(0,m)$,
\[
 0.31184795095143<v_1<0.31184795095144,\qquad
 1.05793812431641<v_2<1.05793812431642,
\]
rounded outward for display. With $v_0=0$ and $v_3=m$, $P$ has constant
sign on each $(v_j,v_{j+1})$. The second integral in \eqref{eq:numS} is
therefore exactly $\sum_{j=0}^{2}|T(v_j)-T(v_{j+1})|$.

\section{Computational certificates}

\subsection*{Two-shift witness}
Three verifiers use the exact decimals in \texttt{witness.json} and
the transform identities in Section 2, and bound the tail by
\eqref{eq:tail}. They use different finite integration routines.

\paragraph{Complex ball quadrature.}
The Python-FLINT verifier uses Arb at 128-bit precision on finite panels.
For $\epsilon=10^{-12}$, it bounds $2\max(y,0)$ by
$y+\sqrt{y^2+\epsilon^2}$ in the numerator correction. For the norm it
integrates $\sqrt{H(z)H(-z)+\epsilon^2}$, an upper bound on the real
axis. Both square roots use Arb's analyticity check on complex boxes.
The resulting outward-rounded inequalities are
\[
 N>0.9177970191817,\qquad D<0.8005327903879 .
\]
These imply $(D/N)^2<0.76079030$.

\paragraph{Trapezoid check.}
The mpmath verifier encloses endpoint values at 35 decimal digits. On a
panel of length $h$, linear interpolation of a real or complex $g$ has
pointwise error at position $s$ at most $s(h-s)\sup|g''|/2$.
Convexity and the Lipschitz constant 1 of modulus and positive part
therefore give a trapezoid upper bound with additive error
$h^3\sup|g''|/12$, even on panels containing a zero. The code bounds
the derivatives term by term. This check gives $(D/N)^2<0.760799446$.

\paragraph{Simpson check.}
A third verifier shares no code with the first two. It isolates the two
sign changes of $P_1$ on $(0,c_1/a)$ by exact rational root isolation and
integrates the positive part in closed form. For the norm it applies
adaptive Simpson panels on $[0,96]$, bounding the fourth derivative of
$|H|$ on each panel by Taylor-ball arithmetic at 150 bits, and uses
\eqref{eq:tail} with $R=96$. This check gives $(D/N)^2<0.760790289$.

\paragraph{Lean verification.}
The accompanying Lean 4 file proves the infinite-tail estimate
\eqref{eq:tail} for finite sums of the indicated type, checks that its right side at $R=64$
is below $0.000002035$ for Table 1, and proves the final
ratio inequality assuming the stated numerical enclosures for $N$ and $D$.
It compiles with Lean and Mathlib v4.24.0 with warnings treated as errors,
and these declarations depend only on the axioms \texttt{propext},
\texttt{Classical.choice} and \texttt{Quot.sound}.
The Python interval verifiers supply the enclosures.

\subsection*{Single-shift witness}
Both verifiers form the rational polynomials exactly, isolate the roots
of $P$ with SymPy, and evaluate the numerator through $T(v)$.

\paragraph{Trapezoid with an exact derivative bound.}
Let $f(\theta)=S(e^{-2\ii\theta})$, $L=\pi/2$, $n=32768$ and $b=L/n$.
Differentiating the trigonometric polynomial twice gives
$\sup|f''|\le K=4\sum_{j=0}^{35} j^2|r_j|=126726777/50000$. The panel
argument of the trapezoid check above, applied to $f$ before taking the
modulus, gives
\begin{equation}
 D_*\le B_n+\frac{KLb^2}{12},\qquad
 B_n=b\left(\frac{|f(0)|+|f(L)|}{2}+\sum_{j=1}^{n-1}|f(jb)|\right).
 \label{eq:trapS}
\end{equation}
With 200-bit Arb balls the verifier obtains, rounded outward,
\[
 N_*>16.68808877572884,\qquad
 B_n<14.55577666588763,\qquad
 \frac{KLb^2}{12}<0.000000762389490,
\]
so $D_*<14.55577742827712$ and $(D/N)^2=(D_*/N_*)^2<0.760777411894$.
This proves the proposition.

\paragraph{Cauchy--Schwarz check.}
The second verifier writes $D_*=\frac12\int_0^\pi|S(e^{\ii\phi})|\,d\phi$
and splits $[0,\pi]$ into 2048 panels $I$. On each panel
$\int_I|S|\le |I|\,(|I|^{-1}\int_I|S|^2)^{1/2}$, and the mean of
$|S|^2$ is a finite trigonometric sum computed exactly from the
autocorrelations of the coefficients. With mpmath interval arithmetic
at 65 digits this gives $(D/N)^2<0.760777857$.

\subsection*{Floating-point cross-checks}
The interval verifiers rely on the reductions in Sections 2 and 3. As a
check on those reductions, a separate floating-point evaluation of the
quotient in \cite[(1.1)]{cmqr} works directly in the variables $t$ and $x$. It compares the closed
form of $F$ with numerical Fourier inversion of $\widehat F$, integrates
$|\widehat F|e^{\pi t}$ between its sign changes and integrates $|F|$ over
the real line. It gives $(D/N)^2\approx0.7607902846$ for the two-shift
witness and $(D/N)^2\approx0.7607773322$ for the single-shift witness,
consistent with the certified upper bounds. These checks are not proofs.

\paragraph{Reproduction and provenance.}
Each of the following commands, run in the archive directory, checks one
of the witnesses:
\begin{itemize}
\item \texttt{python3 verify\_single\_shift\_36\_arb.py} (python-flint 0.9.0, SymPy 1.14.0);
\item \texttt{python3 verify\_single\_shift\_36\_cauchy.py} (mpmath 1.3.0, SymPy 1.14.0);
\item \texttt{python3 check\_arb.py witness.json} (python-flint 0.9.0);
\item \texttt{python3 certify.py witness.json} (mpmath 1.3.0);
\item \texttt{python3 verify\_two\_shift\_simpson.py witness.json}\newline
      (python-flint 0.9.0 and SymPy 1.14.0).
\end{itemize}
The companion archive includes the inputs, these verifiers, the
floating-point cross-checks, the optimization code and the Lean
source\ifReplicationURL{~\cite{replication}}. The two-shift witness was
found by optimizing two translated copies of the test-function basis
of \cite[\S6]{cmqr}. The single-shift witness was found by enlarging a
24-coefficient numerical trial to 36 coefficients at a fixed shift and
dilation and rounding to denominator $10^6$. OpenAI Codex assisted with
the numerical searches, verification code and preparation of this note.
Anthropic's Claude independently reran all verifiers and the Lean check
before deposit.

\begin{thebibliography}{1}
\bibitem{cmqr} E. Carneiro, M. B. Milinovich, E. Quesada-Herrera and
A. P. Ramos, \emph{Fourier optimization, the least quadratic non-residue,
and the least prime in an arithmetic progression}, Mathematics of
Computation, \href{https://doi.org/10.1090/mcom/4154}{doi:10.1090/mcom/4154}.
\href{https://arxiv.org/abs/2404.08380}{arXiv:2404.08380}.
\bibitem{alpoge} L.~Alp\"oge, [Announcement of a Lean formalization of the
Riemann hypothesis], November 20, 2026.
\url{https://alpo.ge/remark.pdf}.

\bibitem{sutherland} A.~V.~Sutherland, \emph{Verification files for the
Lean proof of the Riemann Hypothesis}, MIT, November 20, 2026.
\url{https://math.mit.edu/~drew/grh/}.

\ifReplicationURL{\bibitem{replication} \emph{Replication materials for
Certified Fourier bounds for the least quadratic nonresidue}. Exact
coefficients, verification code, and Lean source.
\expandafter\url\expandafter{\replicationURL}. Source repository:
\url{https://github.com/novalis78/least-nonresidue-two-shift}.}
\end{thebibliography}
\end{document}
